\section{Actor Model vs Account Model}

\subsection{Ethereum and Polkadot: Shared Global State}

Both Ethereum and Polkadot use an \textbf{account model} where smart contracts
operate on shared, globally-accessible state. In Ethereum, any contract can read
the storage of any other contract (via \texttt{SLOAD}). In Polkadot, ink!
contracts are more isolated by default, but the underlying execution environment
is still a shared state machine.

Key characteristics:
\begin{itemize}
    \item \textbf{Reentrancy is possible}: Because the call stack is synchronous,
          a contract can be called again while it is still executing (reentrancy).
    \item \textbf{Global composability}: ContractA can directly call any function
          on ContractB without going through a message queue.
    \item \textbf{MEV (Maximal Extractable Value)}: The shared mempool allows
          validators to reorder transactions, potentially extracting value from
          the chain.
    \item \textbf{Deterministic execution}: All nodes replay the same sequence
          of operations and converge on the same state.
\end{itemize}

\subsection{TON: The Actor Model}

TON implements the \textbf{Actor model}, originally from distributed computing
(Hewitt, Bishop, and Steiger, 1973). Each smart contract is an \textbf{actor}
with:

\begin{itemize}
    \item \textbf{Exclusive private state}: No other contract can read or modify
          a contract's storage directly. State is only accessible through the
          contract's own message handlers.
    \item \textbf{Message-based communication}: Actors communicate exclusively
          via asynchronous messages. There are no direct function calls across
          contract boundaries.
    \item \textbf{No reentrancy}: Since there is no synchronous call stack,
          reentrancy is structurally impossible. A contract cannot be interrupted
          mid-execution by another call.
    \item \textbf{Each actor processes one message at a time}: The TON runtime
          ensures serial message processing per contract, eliminating race
          conditions within a single actor.
\end{itemize}

\subsection{Security Implications for the Chain Pattern}

\begin{table}[h]
\centering
\begin{tabular}{|l|c|c|c|}
\hline
\textbf{Security Property} & \textbf{Ethereum} & \textbf{Polkadot} & \textbf{TON} \\
\hline
Reentrancy risk & High & Medium & None \\
Cross-contract read & Yes & Limited & No \\
Atomic execution & Yes & Yes & No \\
Message ordering & N/A & N/A & Non-deterministic \\
Parallelism & None & None & Shard-level \\
\hline
\end{tabular}
\caption{Security model comparison}
\end{table}

For our A $\to$ B $\to$ C chain pattern:

\begin{itemize}
    \item \textbf{Ethereum}: The chain is atomic and synchronous, but requires
          reentrancy guards (\texttt{ReentrancyGuard} or checks-effects-interactions).
          A malicious ContractB could attempt to re-enter ContractA during the call.
    \item \textbf{Polkadot}: Similar security profile to Ethereum, but ink!'s
          type system and the \texttt{build\_call} API make accidental reentrancy
          less likely.
    \item \textbf{TON}: Reentrancy is impossible, but the chain is non-atomic.
          The professor's grader must check for bounce messages or confirmation
          receipts. If ContractC receives the message but ContractB crashes before
          returning a confirmation, ContractA never knows the chain completed.
\end{itemize}

\subsection{State Isolation and Composability Trade-off}

TON's actor model provides stronger security guarantees (no reentrancy, no
unauthorized state reads) at the cost of composability. In Ethereum, composing
contracts (using the output of one contract as input to another in the same
transaction) is seamless. In TON, such composition requires multi-step
asynchronous protocols, which are more complex to design and verify.

For educational purposes, this trade-off is instructive: students should
understand that stronger isolation means weaker composition, and that the choice
of platform determines which properties are guaranteed by the runtime versus
enforced by the developer.